% Module: two-moment theorem (proof), counting lemma, proofs of corollaries. Labels: lem:count
\begin{proof}[Proof of Theorem~\ref{thm:two}]
\emph{(a) Reduction.} $J_I$ and $T_I$ are equivariant under $Z\mapsto UZV^\top$, $U\in O(K)$, $V\in O(d)$, because
$\polar(UZV^\top)=U\polar(Z)V^\top$ and the law of $Z$ is invariant. $E=\R^K\otimes\R^d$ is irreducible, so
$J_I=c\,\mathrm{Id}$ with $c=\tr J_I/N=\E\inner{\polar(Z)}Z/N=c_{K,d}$ \citep{xie2026muon,chen2026when}. $T_I$ is an invariant element of
$E^{\otimes4}=(\R^K)^{\otimes4}\otimes(\R^d)^{\otimes4}$; the invariants of $O(K)\times O(d)$ are spanned by the
products $\pi\otimes\pi'$ of an $O(K)$-invariant and an $O(d)$-invariant, and by the first fundamental theorem for the
orthogonal group each factor is spanned by the three pairings of the four indices. Label a pairing by the input slot
paired with the output index; $T_I$ is symmetric in its three input slots, the symmetric group $S_3$ permutes the nine
products $(\pi,\pi')$ with two orbits, $\{\pi=\pi'\}$ and $\{\pi\ne\pi'\}$, and the orbit averages are the tensors of
$\cA$ and $\cB$ on symmetric arguments. Hence $T_I=\alpha\cA+\beta\cB$. To find $\alpha,\beta$ contract with one representative of each orbit:
using $\He_3(Z)^{(ab),(cd),(ef)}=Z_{ab}Z_{cd}Z_{ef}-Z_{ab}\delta_{(cd),(ef)}-Z_{cd}\delta_{(ab),(ef)}-Z_{ef}\delta_{(ab),(cd)}$,
\begin{align*}
\textstyle\sum_{ij}\cA\big[\E\,\polar(Z)_{ij}\He_3(Z)\big]_{ij}&=\E\big[\nnuc Z(\nF Z^2-N-2)\big]=-m_1,\\
\textstyle\sum_{ij}\cB\big[\E\,\polar(Z)_{ij}\He_3(Z)\big]_{ij}&=\E\big[\inner{\polar(Z)}{ZZ^\top Z}\big]-(K+d+1)m_1=m_3-(K+d+1)m_1,
\end{align*}
where $\E[\nnuc Z\nF Z^2]=(N+1)m_1$ because $\nF Z$ is independent of $Z/\nF Z$ and $\E\chi_N^3=(N+1)\E\chi_N$. On the
other side, by Lemma~\ref{lem:count}, the same contractions of $\alpha\cA+\beta\cB$ are
$\frac N3[(N+2)\alpha+(K+d+1)\beta]$ and $\frac N3[(K+d+1)\alpha+\frac12(N+K+d+3)\beta]$. The determinant of this system
is $\frac{N^2}{18}(K-1)(d-1)(K+2)(d+2)$; solving, and substituting $m_3=(K+d-q)m_1$ from (b), gives
$\alpha=-c_{K,d}a_1$, $\beta=-c_{K,d}a_2$ with \eqref{eq:a12}. Then \eqref{eq:two} is Theorem~\ref{thm:cf} with the scaling
of Lemma~\ref{lem:gauss}. The map $K_3\mapsto(\cA[K_3],\cB[K_3])$ is onto $E\oplus E$ (the two equivariant maps are
linearly independent and $E$ is irreducible), so $K_3\mapsto a_1\cA+a_2\cB$ is onto $E$.

\emph{(b)} $\sum_is_i^3=\inner Z{f(Z)}$ with the spectral function $f(Z)=(ZZ^\top)^{1/2}Z=U\operatorname{diag}(s_i^2)V^\top$,
which is $C^1$; by Stein's identity $m_3=\E\operatorname{div}f(Z)$ and by the divergence formula for spectral functions
\citep{candes}, with $\varphi(s)=s^2$,
$\operatorname{div}f=\sum_i\big(\varphi'(s_i)+|K-d|\frac{\varphi(s_i)}{s_i}\big)+\sum_{i\ne j}\frac{s_i\varphi(s_i)-s_j\varphi(s_j)}{s_i^2-s_j^2}
=(K+d)\sum_is_i-\sum_{i\ne j}\frac{s_is_j}{s_i+s_j}$, since
$\frac{s_i^3-s_j^3}{s_i^2-s_j^2}=s_i+s_j-\frac{s_is_j}{s_i+s_j}$. The bounds on $q$ follow from
$0<\frac{s_is_j}{s_i+s_j}<\frac{s_i+s_j}4$ almost surely, and $a_2>0$ from \eqref{eq:a12}.

\emph{(c)} Algebra from \eqref{eq:a12}, using $Kd+1-(K+d)=(K-1)(d-1)$ and, for $K\le d$,
$Kd+2-K(K+d+1)=-(K-1)(K+2)$; positivity from $q<\frac{r-1}2\le\frac{\bar d-1}2$.

\emph{(d)} $\nabla^4\psi_I(0)[D,D,D,D]=\inner D{T_I[D^{\otimes3}]}$, $\cA[D^{\otimes3}]=\nF D^2D$,
$\cB[D^{\otimes3}]=DD^\top D$. If $a_1\ge0$ the sign is clear. If $a_1<0$, $\nF D^4\le\operatorname{rank}(D)\tr(DD^\top)^2$
gives $a_1\nF D^4+a_2\tr(DD^\top)^2\ge(r\,a_1+a_2)\tr(DD^\top)^2>0$ by (c).

\emph{(e)} The empirical distribution of the $s_i^2/d$ converges almost surely to the Marchenko--Pastur law of ratio
$\gamma$, with convergence of the moments $m_1/(K\sqrt d)$ and $m_3/(Kd^{3/2})$ (the extreme singular values are bounded in
every $L^p$); by (b), $q/K=(K+d)/K-m_3/(Km_1)\to[(1+\gamma)\mu_{1/2}-\mu_{3/2}]/(\gamma\mu_{1/2})$. The hypergeometric
form of $\mu_s$ is the Narayana expansion for integer $s$; for $s=\frac12,\frac32$ only the integral definition of $\mu_s$ is used; $\mu_{1/2}(1)=\frac8{3\pi}$, $\mu_{3/2}(1)=\frac{64}{15\pi}$
(quarter-circle law), hence $\kappa(1)=2-\frac85$. The limits of $a_1,a_2$ follow from \eqref{eq:a12}.

\emph{$K=1$.} The same computation with a single invariant; $m_3=(d+1)m_1$ and the contraction of the invariant with
itself is $N(N+2)/3$, so $\alpha+\beta=-3m_1/(d(d+2))$.
\end{proof}

\begin{lemma}[counting]\label{lem:count}
Let $\mathsf A_1$, $\mathsf B_1$ be one pairing tensor from each orbit and $\bar{\mathsf A}$, $\bar{\mathsf B}$ the orbit
averages (so that $\bar{\mathsf A}$, $\bar{\mathsf B}$ are the tensors of $\cA$, $\cB$ on symmetric arguments). Then
$\inner{\mathsf A_1}{\bar{\mathsf A}}=\frac{N(N+2)}3$, $\inner{\mathsf A_1}{\bar{\mathsf B}}=\inner{\mathsf B_1}{\bar{\mathsf A}}=\frac{N(K+d+1)}3$,
$\inner{\mathsf B_1}{\bar{\mathsf B}}=\frac{N(N+K+d+3)}6$.
\end{lemma}
\begin{proof}
The inner product of two pairing tensors is $K^{\ell}d^{\ell'}$, $\ell,\ell'\in\{1,2\}$ being the numbers of closed
loops of row and of column indices. The three inner products of $\mathsf A_1$ with the first orbit are $N^2,N,N$; its
six inner products with the second orbit are $Kd^2,Kd^2,K^2d,K^2d,Kd,Kd$; the six inner products of $\mathsf B_1$ with
its orbit are $K^2d^2,K^2d,Kd^2,Kd,Kd,Kd$. (Direct enumeration; also verified by brute-force contraction for seven sizes.)
\end{proof}

\begin{proof}[Proof of Corollary~\ref{cor:dim}]
For independent entries the only nonzero components of $K_3$ are $K_3^{(ij),(ij),(ij)}=\kappa_3$, so
$\cA[K_3]_{ij}=\cB[K_3]_{ij}=\kappa_3$. The value of $a_1+a_2$ is Theorem~\ref{thm:two}(c); the upper bound is
$q_{K,d}<\frac{r-1}2$ with $(r+2)(\bar d+2)=(K+2)(d+2)$, and the lower bound is Theorem~\ref{thm:a1}.
\end{proof}

\begin{proof}[Proof of Corollary~\ref{cor:rank1}]
If $\zeta=uv^\top$ then $\zeta\zeta^\top\zeta=|u|^2|v|^2uv^\top=\nF\zeta^2\zeta$, so $\cB=\cA$ and
$\gs=(a_1+a_2)\cA/(6\sigma^2)$. For the Frobenius norm apply the case $K=1$ of Theorem~\ref{thm:two} on $\R^{Kd}$:
$\gs=\frac3{Kd+2}\cA/(6\sigma^2)$. The ratio is $\varrho=\frac{(3+2q)(Kd+2)}{(K+2)(d+2)}$; it exceeds
$\frac{Kd+2}{K+d+1}>1$ by Theorem~\ref{thm:a1} (as $Kd+2-(K+d+1)=(K-1)(d-1)>0$), and it is below
$\frac{(r+2)(Kd+2)}{(K+2)(d+2)}=\frac{Kd+2}{\bar d+2}<r$.
\end{proof}

\begin{proof}[Proof of Corollary~\ref{cor:compare}]
(i) Theorems~\ref{thm:cf} and~\ref{thm:two}; for the $\ell_1$ norm and diagonal $\Sigma$,
$J=\operatorname{diag}(\sqrt{2/\pi}/\sigma_k)$ and $T[K_3]_k=-\sqrt{2/\pi}\,\kappa_{3,k}/\sigma_k^3$, so
$(\gs)_k=\kappa_{3,k}/(6\sigma_k^2)$. (ii) The only nonzero entries of $K_3$ are the permutations of
$((11),(12),(12))$, whence $\cA=\cB=\tau e_1e_1^\top$. Such a law exists:
$\zeta_{11}=\varrho(\zeta_{12}^2-1)/\sqrt2+\sqrt{1-\varrho^2}\,\epsilon$ with $\epsilon$ an independent centred Gamma
variable whose skewness cancels that of the first term. (iii) For this tensor
$\cA_{11}=1+K_3^{(11),(22),(22)}$, $\cB_{11}=1$ and all other entries of $\cA,\cB$ vanish (the pattern
$(kj),(kl),(il)$ of $\cB$ is never a permutation of $((11),(22),(22))$), so
$a_1\cA+a_2\cB=(a_1+a_2+a_1K_3^{(11),(22),(22)})e_1e_1^\top=0$; a small multiple of this $K_3$ is the third cumulant of a law
with covariance $I$ and a density (the pairs of covariance and third moment of laws with a density near the Gaussian
fill a neighbourhood of $(I,0)$). (iv) Cumulants add over independent summands, and
$a_1\nF M^4+a_2\tr(MM^\top)^2>0$ by Theorem~\ref{thm:two}(d).
\end{proof}
